\documentclass[titlepage]{ltjsarticle}
\usepackage{url}
\usepackage{listings}

\title{計算機科学基礎実験 レポート2}
\author{氏名：加藤悠菜 / 学籍番号：6325040}
\date{\today}

\begin{document}

\maketitle

\section{課題1:2進数の個数}
\subsection{実装}
\subsubsection{アルゴリズムの説明}
与えられた$n$を2進数で表記した際の1の個数を求めるため、再帰関数 \texttt{count\_ones} を実装した。
このアルゴリズムは、$n \pmod 2 $によって最下位ビットが1であるか確認し、その後$n/2$を再帰的に呼び出すことで、全ての桁を調べることを可能とした。

\subsubsection{プログラム}
\begin{verbatim}
let rec count_ones n=
if n<0 then failwith "Error"
else
match n with
0 -> 0
|_ -> (n mod 2) + count_ones (n / 2)
;;
\end{verbatim}

\subsection{実行過程}
私の学籍番号の下二桁である$n=40$を例に実行過程を説明する。
\begin{enumerate}
    \item \texttt{count\_ones 40} $\rightarrow$ $40 \pmod 2 = 0$を保持し、\texttt{count\_ones 20}を呼び出し。
    \item \texttt{count\_ones 20} $\rightarrow$ $20 \pmod 2 = 0$を保持し、\texttt{count\_ones 10}を呼び出し。
    \item \texttt{count\_ones 10} $\rightarrow$ $10 \pmod 2 = 0$を保持し、\texttt{count\_ones 5}を呼び出し。
    \item \texttt{count\_ones 5} $\rightarrow$ $5 \pmod 2 = 1$を保持し、\texttt{count\_ones 2}を呼び出し。
    \item \texttt{count\_ones 2} $\rightarrow$ $2 \pmod 2 = 0$を保持し、\texttt{count\_ones 1}を呼び出し。
    \item \texttt{count\_ones 1} $\rightarrow$ $40 \pmod 2 = 1$を保持し、\texttt{count\_ones 0}を呼び出し。
    \item \texttt{count\_ones 0} $\rightarrow$ $n=0$より返り値$0$を返し、合計は$0+0+0+1+0+1+0=2$となる。
\end{enumerate}
商は「次の桁への移動」、余りは「現在の桁の1の判定」という役割を担っている。

\section{課題2:べき乗$n^n$の計算}
\subsection{$n^n$の値を求める}
\subsubsection{アルゴリズムの説明}
このプログラムでは正の整数$n$を引数に取り、$n$の$n$乗を計算する。
$n^n$は「$n$を$n$回かける」という操作であるため、再帰を用いる場合は「現在何回かけたか」を管理するカウンターが必要となる。
しかし、引数は1つしかないため、関数内部に補助関数(\texttt{aux})を定義した。

\subsubsection{プログラム}
\begin{verbatim}
    let rec power_val n=
 let rec aux count=
  if count = 0 then 1
  else n * aux (count - 1)
in
aux n
;;
\end{verbatim}

\subsection{呼び出し回数を数える}
\subsubsection{アルゴリズムの説明}
\texttt{powwer\_val}で実装した$n$の乗算ではなく、「1を足し合わせる」という操作に置き換わる。
再帰の各ステップで\texttt{1 + aux(count - 1)}を行うことで、関数の呼び出し回数をカウントしている。

\subsubsection{プログラム}
\begin{verbatim}
    let rec power_steps n=
 let rec aux count=
  if count = 0 then 1
  else 1 + aux (count - 1)
 in
 aux n
;;
\end{verbatim}

\section{課題3:コラッツ予想}
コラッツ予想に基づき、任意の正の整数$n$が1になるまでの操作を実装した。
基本規則は以下である。
\begin{itemize}
    \item $n$が偶数の場合:$n\leftarrow n/2$
    \item $n$が奇数の場合:$n\leftarrow 3n+1$
\end{itemize}
\subsection{ステップ回数を求める}
\subsubsection{アルゴリズムの説明}
関数\texttt{collatz\_steps}は1になるまでの操作回数を求める。
再帰の各ステップで、$n=1$でなければ操作回数として「1」を足し、規則に基づいた次の値を引数として再帰呼び出しを行う。
これにより、$n=1$までの操作回数がカウントされる。

\subsubsection{プログラム}
\begin{verbatim}
let rec collatz_steps n=
if n=1 then 0
else if n mod 2 = 0 then
 1 + collatz_steps (n/2)
else
 1 + collatz_steps (3 * n + 1)
;;
\end{verbatim}

\subsection{推移のリストを求める}
\subsubsection{アルゴリズムの説明}
関数\texttt{collatz\_path}は計算過程の数値を順に並べたリストを求める。
リスト構成子(\texttt{::})を使い、現在の値$n$をリストの先頭に追加しながら再帰呼び出しを行う。
これにより、初項から1に至るまでの過程が保存される。

\subsubsection{プログラム}
\begin{verbatim}
let rec collatz_path n=
if n=1 then [1]
else if n mod 2 =0 then
 n :: collatz_path (n/2)
else
 n :: collatz_path (3 * n + 1)
;;
\end{verbatim}

\section{課題4:べき乗とコラッツ予想に関する考察}
\subsection{「入力値の大きさ」と「ステップ数」の関係:べき乗vsコラッツ予想}
$n=7, 8, 9, 10$ の範囲でそれぞれのステップ数を調べた。

\begin{table}[h]
\centering
\begin{tabular}{|c|c|c|}
\hline
入力 $n$ & べき乗（power\_steps） & コラッツ予想（collatz\_steps） \\ \hline
7 & 8 & 16 \\ \hline
8 & 9 & 3 \\ \hline
9 & 10 & 19 \\ \hline
10 & 11 & 6 \\ \hline
\end{tabular}
\caption{入力値とステップ数の比較}
\end{table}

\clearpage
この結果から、ステップ数の性質には違いが見られる。
「べき乗のアルゴリズム」においては、ステップ数は常に $n+1$ であり、入力値に対して比例的に増加する。これは動作が完全に予測可能であることを示している。

対して、「コラッツ予想」のステップ数は入力に対して極めて不規則である。例えば、$n=8$ は $n=7$ よりも大きい入力であるにも関わらず、ステップ数は16から3へと大幅に減少している。これは $8 = 2^3$ であり、「偶数なら2で割る」というルールのみが適用され続けるためである。
このように、コラッツ予想においては入力の大きさよりも、その数値が「2のべき乗に近いか」や「奇数ルール（$3n+1$）によって数値がどれだけ増大するか」がステップ数に影響を与えることが分かった。

\subsection{数値の「合流」と収束の仕組み}
\texttt{collatz\_path} を実行した際の結果を比較すると、異なる初期値からスタートした数値が、計算の途中で同じ数値に到達し、それ以降全く同一の経路を通っていた。
これが「合流」である。

例えば、$n=9$ の推移を見ると、途中で $10$ に到達している。
\begin{itemize}
    \item \texttt{collatz\_path 9} $\rightarrow$ [9; 28; 14; 7; 22; 11; 34; 17; 52; 26; 13; 40; 20; \textbf{10}; 5; 16; 8; 4; 2; 1]
    \item \texttt{collatz\_path 10} $\rightarrow$ [\textbf{10}; 5; 16; 8; 4; 2; 1]
\end{itemize}
この時点で $n=9$ のルートは $n=10$ のルートに合流したと言える。
また、$n=7$ や $n=3$ なども、最終的には $16 \to 8 \to 4 \to 2 \to 1$ という「2のべき乗のサイクル」に吸い込まれる形で収束していく。
このように1度でも経験済みの数値や2のべき乗に合流すれば最終的に数値は1に収束することがわかった。
\end{document}